(1) 定义惯量矩$I_G=\sum_im_ir_i^2$，两次时间微分给$\frac12\ddot I_G=2K+\sum_i\boldsymbol r_i\cdot\boldsymbol F_i$。对于Newton引力势$U\propto r^{-1}$，Euler齐次关系$\sum_i\boldsymbol r_i\cdot\boldsymbol F_i=U$。处于长期统计稳态且$\langle\ddot I_G\rangle=0$时得到$2\langle K\rangle+\langle U\rangle=0$；这不是任一瞬时位置都严格满足的恒等式。
(2) 稳态比能量$E=\langle K+U\rangle=-\langle K\rangle$。按动能定义有效温度$\langle K\rangle=(3/2)Nk_BT$，从而$E=-(3/2)Nk_BT$，对该族准稳态孤立平衡分支的微正则导数$C=(\partial E/\partial T)_N=-3Nk_B/2<0$。负热容允许于长程自引力、非加性系统，不代表粒子单体能量为负。
(3) 若星团通过辐射或逃逸缓慢损失总能量$\Delta E<0$并重新Virial化，则$\Delta T=-2\Delta E/(3Nk_B)>0$；引力势能变得更负，释放出的势能超过向外输送的能量，剩余转成随机动能。与普通正热容样品不同，失能反而升温可导致核心收缩与gravothermal失稳。
(4) 表面压力不可忽略时Virial定理应加$-3PV$型边界项（符号依定义），逃逸会改变$N$并携走能量，高密核心则受到碰撞、双星和量子效应影响。负微正则热容不能作为与无限温度热库稳定正则接触的平衡支：正则系综有$C_V=\operatorname{Var}(E)/(k_BT^2)\ge0$，且无限自引力点粒子气体无短程正规化时配分函数可能发散。
